\subsection*{Declaração das funções}

\begin{verbatim}
//definição da func. rosenbrock segundo o arquivo enviado
function [x] = funcR(a,b)
x = 100*(0.004*b-(0.004*a)^2)^2+(1-0.004*a)^2
endfunction

//definição da função z
function [x] = funcZ(a,b)
x = a*sin(sqrt(abs(a))) - b*sin(sqrt(abs((b))))
endfunction

//definição das funções pertinentes à w13
function [x] = funcw1(a, b)
a1 = funcR(a, b)
a2 = funcZ(a, b)
x = a1+a2
endfunction

function [x] = funcw3(a, b)
a1 = funcR(a,b)
a2 = funcZ(a,b)
x = a1-a2
endfunction

function [x] = funcw4ant(a,b)
a1 = funcR(a,b)
a2 = funcZ(a,b)
x = sqrt(a1^2 + a2^2)
endfunction

//definição de w13, func. a ser minimizada
function [x] = funcw13(a,b)
x = sqrt(abs(funcw1(a,b))) + sqrt(abs(funcw3(a,b))) - funcw4ant(a,b)*cos(0.004*a)
endfunction

\end{verbatim}

\subsection*{Algoritmo da Nuvem de Partículas}
\begin{verbatim}

//inicialização do vetor de população de tamanho m=101, localizados numa diagonal do espaço de busca
xpop = [-500:10:500;-500:10:500]
//inicialização do vetor de melhor posição alcançada para um membro da pop.
ppop = [-500:10:500;-500:10:500]
//inicialização do vetor velocidade para cada membro da população
vpop = zeros(2, 101)
//inicialização das variáveis de melhor/mínimo global, g1 corresponde ao eixo x e g2 ao eixo y
g1 = 0
g2 = 0 
//constantes pertinentes ao método usado (particle swarm, enxame de partículas) 
c1 = 2
c2 = 2
peso = 0.70

for k=1:50
// condição de término é a de rodar 100 iterações do código
    for l=1:101
//para cada um dos membros da população definida rode:
        if funcw13(xpop(1,l),xpop(2,l))<funcw13(ppop(1,l),ppop(2,l)) then
           ppop(1,l) = xpop(1,l)
           ppop(2,l) = xpop(2,l)
           //se o valor de w13 é menor para esta posição xpop que ppop, grave em ppop
        end
           if funcw13(xpop(1,l),xpop(2,l))<funcw13(g1,g2) then
           g1 = xpop(1,l)
           g2 = xpop(2,l)
           //se o valor de w13 é menor para esta posição de xpop que g1 e g2, grave em g1 e g2
       end
        r11 = rand()
        r12 = rand()
        r21 = rand()
        r22 = rand()
//escolha dos números randômicos a cada iteração do laço
        vpop(1,l) = peso*vpop(1,l) + c1*r11*(ppop(1,l) - xpop(1,l)) + c2*r12*(g1 -         xpop(1,l))
        vpop(2,l) = peso*vpop(2,l) + c1*r21*(ppop(2,l) - xpop(2,l)) + c2*r22*(g2 -         xpop(2,l))
//determinação das velocidades
        xpop(1,l) = xpop(1,l) + vpop(1,l)
        xpop(2,l) = xpop(2,l) + vpop(2,l)
//determinação das posições
        if xpop(1,l)>=500 then
            xpop(1,l) = 500
        end
        if xpop(1,l)<=-500 then
            xpop(1,l) = -500
        end
        if xpop(2,l)>=500 then
            xpop(2,l) = 500
        end
        if xpop(2,l)<=-500 then
            xpop(2,l) = -500
        end
//projeção do espaço de busca nos eixos x e y (nos limitamos ao retângulo delimitado por -500<=x<=500 e -500<=y<=500.)
    end
end
\end{verbatim}

\subsection*{Algoritmo Genético e funções pertinentes à cadeias de bit}
\begin{verbatim}

//retornar a cadeia de bits para o membro da população index em xpop ou ypop
function [x] = extCBit(a, index)
    s = size(a)
    x = zeros(s(1), 1)
    for h=1:s(1)
        x(h, 1) = a(h, index)
    end
endfunction

// converter o valor da cadeia de bits para o valor real dados os limites
function [x] = retCBit(a, LimS, LimI)
    x = 0
    for h=1:length(a)
        x = x + a(h)*2^(h-1)
    end
    x = LimI+x*(LimS - LimI)/(2^(length(a)+1)-1)
endfunction

//função de randomização de um bit
function [x] = ranBit(bit)
x = bit
    if rand()>0.5 then
        x= 1
    else
        x=0       
    end
endfunction

//função de randomização de uma cadeia de bits
function [a] = ranCBit(a)
    for h=1:length(a)
        a(h)= ranBit(a(h))
    end
endfunction

//cálculo da aptidão de duas cadeias de bit x e y
function [x] = aptCBit(a, b)
    a1 = retCBit(a, 500, -500)
    a2 = retCBit(b, 500, -500)
    x = 1/(funcw13(a1, a2) + 600)
endfunction

//função crossover entre duas cadeias de bit
function [x, y] = crossover(a,b)
    l = rand()
    pcorte = (l - modulo(l, 1/length(a)))*length(a)
    a1 = a
    a2 = b
    for h=1:pcorte
        a(h) = a2(h)
        b(h) = a1(h)
    end
    x = a
    y = b
endfunction

//função mutação para cada bit
function [x] = mutacao(a,mut)
    x=a
    if rand()<mut then
        if a==0 then
            x=1
        else
            x=0
        end
    end
endfunction

function [x] = mutCBit(a)
    a1 = length(a)
    x = zeros(a1, 1)
    for h=1:length(a)
    x(h, 1) = mutacao(a(h, 1), mut)
    end
endfunction



// N é o numero de bits nas cadeias de bits, 10 bits para x e 10 bits para y
s = size(xpop)
m = s(1)
N = s(2)
//variáveis auxiliares
cbitANx = zeros(16, 1)
cbitANy = zeros(16, 1)
ubitANx = zeros(16, 1)
ubitANy = zeros(16, 1)
//inicialização melhor par ordenado de cadeias
mbitAPTx=zeros(16,1)
mbitAPTy=zeros(16,1)
//aptidão total
VarAPTt = 0
//chance de seleção para cada membro da população
p = zeros(1,20)
mut = 0.01
cros = 0.75
//inicialização dos membros da população
xpop = zeros(16,20)
ypop = zeros(16,20)
for h=1:16
    for e=1:20
    xpop(h,e) = ranBit(xpop(h,e))
    ypop(h,e) = ranBit(ypop(h,e))
    end
end


//início do algoritmo
// k é o número de iterações do algoritmo
for k=1:36
    for g=1:m
        //cálculo do total de todas as aptidões
        cbitANx = extCBit(xpop, g)
        cbitANy = extCBit(ypop, g)
        VarAPTt = VarAPTt + aptCBit(cbitANx, cbitANy)
        if aptCBit(cbitANx, cbitANy)>aptCBit(mbitAPTx, mbitAPTy) then
            mbitAPTx = cbitANx
            mbitAPTy = cbitANy
        end
        end
    end
     //cálculo das probabilidades de seleção de cada membro
    for g=1:m
        p(1, g) = aptCBit(cbitANx, cbitANy)/VarAPTt
    end
    
    for g=1:m
        if rand()<cros then
            for u=1:m
                if rand()<p(1, u) then
                    ubitANx = extCBit(xpop, u)
                    ubitANy = extCBit(ypop, u)
                    [ubitANx, cbitANx] = crossover(ubitANx, cbitANx)
                    [ubitANy, cbitANy] = crossover(ubitANy, cbitANy)
                    for h=1:N
                    xpop(g, h) = cbitANx(g, 1)
                    ypop(g, h) = cbitANy(g, 1)
                    xpop(u, h) = ubitANx(u, 1)
                    ypop(u, h) = ubitANy(u, 1)                   
                    end
                end
            end
        end
       //mutação das cadeias de bit em x e y para o membro g da população
       for g=1:m
        cbitANx = extCBit(xpop, g)
        cbitANy = extCBit(ypop, g)
        cbitANx = mutCBit(cbitANx)
        cbitANy = mutCBit(cbitANy)
        for h=1:N
            xpop(g, h) = cbitANx(g, 1)
            ypop(g, h) = cbitANy(g, 1)
        end    
       end
end
\end{verbatim}